\section{Experiments}
\label{sec:experiments}

\subsection{Setup and Evaluation}
\label{subsec:experimental_setup}
\label{subsec:experimental_metrics}

The platform uses a 145\,g quadrotor and a 0.9525\,m cable
weighing 17\,g with tracking markers. OptiTrack measures vehicle
pose and cable motion. The vehicle starts at $(0,0,1.4)$\,m,
with target $(1.25,0,1.25)$\,m and strike direction $+x$.
Commands are executed at 30\,Hz without online cable feedback.
The predictor uses 12 cable nodes and 150\,Hz rollout samples;
the command spline has 12 control points over $T=1.5$\,s.
Correction uses $(w_{\mathrm{tip}},w_q,w_d)=(1,0.1,0.01)$
in~\eqref{eq:reference_tracking_cost}.

Five flights were collected for each model--command pair:
the initial M0 plan and two corrections, M1 and M2, toward
the same M0 trajectory and timing. In each preceding batch,
recordings 1, 2, and 4 contribute to fitting; recordings 3
and 5 support development checks. M2 retains the M1 vehicle
network and updates the nominal vehicle response and cable
model. It is selected using the two M1 development recordings
and frozen before the five M2 flights.

Tip and vehicle tracking errors, $E_{\mathrm{tip}}$ and $E_q$,
are three-dimensional RMSEs against the original M0 reference
at matching timestamps over $[0,T_h]$, with $T_h=34/30$\,s.
Target metrics are $d_{\min}$ from~\eqref{eq:target_distance}
over $[0,1.5]$\,s and distance $d_\star$ at the original
strike time $t_\star=1.11725$\,s. Tip speed is estimated from
local position derivatives. Clock alignment uses measured
vehicle streams; missing observations are masked. Statistics
average per-flight metrics, without spatial or target-based
temporal alignment.

\subsection{Results}
\label{subsec:physical_results}
\label{subsec:prediction_results}

\begin{table}[t]
\centering
\caption{Physical performance: mean $\pm$ sample standard
deviation, five flights per batch. Distances are in cm;
tip speed at $t_\star$ is in m/s.}
\label{tab:physical_performance}
\small
\setlength{\tabcolsep}{3pt}
\begin{tabular}{@{}lccc@{}}
\toprule
Metric & M0 & M1 & M2 \\
\midrule
$d_\star$ & $24.12\pm6.94$ & $11.18\pm4.04$ & $13.10\pm7.52$ \\
$d_{\min}$ & $14.73\pm6.02$ & $7.83\pm1.82$ & $8.60\pm6.55$ \\
$E_{\mathrm{tip}}$ & $18.09\pm1.80$ & $16.06\pm2.34$ & $13.28\pm2.60$ \\
$E_q$ & $15.66\pm3.58$ & $11.17\pm2.87$ & $12.63\pm5.65$ \\
Tip speed & $5.98\pm0.53$ & $5.85\pm0.32$ & $5.82\pm0.79$ \\
\bottomrule
\end{tabular}
\end{table}

Table~\ref{tab:physical_performance} shows that M1 substantially
improves target approach. M2 further reduces tip-reference
RMSE by 17.4\% relative to M1, but target accuracy and
repeatability worsen. Mean tip speed remains similar.
Thus, better average trajectory tracking does not ensure
better target approach.

An exploratory subset selected after reviewing the outcomes,
M2 recordings 2 and 4, gives mean $d_\star$, $d_{\min}$,
and $E_{\mathrm{tip}}$ of 6.17, 4.41, and 10.57\,cm.
No independent fault evidence supports excluding the other
three flights; the primary comparison retains all five.

On the two M1 recordings used for model selection, M2 reduces
vehicle-prediction RMSE from 7.46 to 6.72\,cm and command-driven
tip-prediction RMSE from 9.95 to 8.46\,cm. These are development
comparisons using the same commands and measured initial
states, not independent test results.

The study uses successive batches at one target. Small cable
motions observed during hover differ from the stationary
planning state; their effect, including possible propeller
airflow, remains unquantified. Fixed-time error is also
sensitive to clock alignment: a $\pm10$\,ms shift changes
M2 recording 2's $d_\star$ from 2.64 to 6.80--7.15\,cm.

